% Vietnamese translation for OI Public Library
% Source-PDF: IOI中国国家候选队论文/2009/罗穗骞/后缀数组——处理字符串的有力工具.pdf
\documentclass[11pt]{article}
\usepackage{graphicx}
\usepackage{oipl-vn}

\newcommand{\Oh}{\mathcal{O}}
\newcommand{\SA}{\operatorname{SA}}
\newcommand{\Rank}{\operatorname{Rank}}
\newcommand{\Height}{\operatorname{height}}
\newcommand{\Suffix}{\operatorname{suffix}}
\newcommand{\LCP}{\operatorname{LCP}}
\newcommand{\RMQ}{\operatorname{RMQ}}
\renewcommand{\figurename}{Hình}
\newcommand{\luosafigure}[3]{%
  \begin{figure}[htbp]
  \centering
  \raisebox{-\height}{\scalebox{1}[-1]{\includegraphics[width=#2\linewidth]{assets/ioi/2009/luo-suiqian-suffix-array/figure-#1.png}}}
  \caption{#3}
  \end{figure}
}

\title{Mảng hậu tố: công cụ mạnh để xử lý xâu}
\author{Luo Suiqian\\Trường Trung học Phụ thuộc Đại học Sư phạm Hoa Nam}
\date{Luận văn đội tuyển dự tuyển quốc gia IOI 2009}

\begin{document}
\maketitle

\origtitle{Mảng hậu tố: công cụ mạnh để xử lý xâu}
\sourcepath{IOI China candidate team papers / 2009 / Luo Suiqian}

\begin{abstract}
Mảng hậu tố là một công cụ rất mạnh trong xử lý xâu. Nó là một thay thế tinh
tế cho cây hậu tố: dễ cài đặt hơn, có thể thực hiện nhiều chức năng tương tự
với độ phức tạp thời gian không kém nhiều, và dùng ít bộ nhớ hơn đáng kể. Do
đó trong thi Tin học, mảng hậu tố thường thực dụng hơn cây hậu tố.

Bài viết gồm hai phần. Phần thứ nhất giới thiệu hai cách xây dựng mảng hậu tố:
thuật toán nhân đôi và thuật toán DC3, nhấn mạnh cách viết mã ngắn gọn, hiệu
quả, rồi so sánh hai thuật toán. Phần thứ hai trình bày các ứng dụng thường
gặp của mảng hậu tố trong những lớp bài toán xử lý xâu.
\end{abstract}

\paragraph{Từ khóa.}
Xâu; hậu tố; mảng hậu tố; mảng hạng; sắp xếp cơ số; DC3; tiền tố chung dài
nhất.

\tableofcontents

\section{Cài đặt mảng hậu tố}

Phần này trình bày hai phương pháp xây dựng mảng hậu tố: thuật toán nhân đôi
và thuật toán DC3. Cả hai đôi khi bị xem là khó hiểu và khó cài đặt. Vì vậy
ngoài ý tưởng, phần này giữ lại các đoạn mã C ngắn gọn trong bài gốc: bản
nhân đôi khoảng 25 dòng, còn bản DC3 khoảng 40 dòng.

\subsection{Các định nghĩa cơ bản}

\paragraph{Xâu con.}
Với xâu \(r\), xâu con \(r[i..j]\), \(i\le j\), là đoạn từ vị trí \(i\) tới
vị trí \(j\), tức xâu tạo bởi
\[
  r[i],r[i+1],\ldots,r[j].
\]

\paragraph{Hậu tố.}
Hậu tố là xâu con đặc biệt bắt đầu tại một vị trí \(i\) và kéo dài tới cuối
xâu. Hậu tố của \(r\) bắt đầu ở ký tự thứ \(i\) được viết
\[
  \Suffix(i)=r[i..\operatorname{len}(r)].
\]

\paragraph{So sánh từ điển.}
So sánh hai xâu \(u,v\) theo thứ tự từ điển nghĩa là lần lượt so sánh
\(u[i]\) và \(v[i]\) từ \(i=1\). Nếu hai ký tự bằng nhau thì tăng \(i\), nếu
\(u[i]<v[i]\) thì \(u<v\), còn nếu \(u[i]>v[i]\) thì \(u>v\). Nếu một trong
hai xâu đã hết mà vẫn chưa phân định được, xâu ngắn hơn được coi là nhỏ hơn;
nếu hai độ dài bằng nhau thì hai xâu bằng nhau.

Từ định nghĩa này, hai hậu tố của cùng một xâu nhưng bắt đầu ở hai vị trí khác
nhau không thể bằng nhau, vì điều kiện cần để hai xâu bằng nhau là chúng có
cùng độ dài.

\paragraph{Mảng hậu tố.}
Cho xâu \(S\) dài \(n\). Mảng hậu tố \(\SA\) là một hoán vị của
\(1,2,\ldots,n\) thỏa mãn
\[
  \Suffix(\SA[i]) < \Suffix(\SA[i+1]),\qquad 1\le i<n.
\]
Nói cách khác, ta sắp xếp mọi hậu tố theo thứ tự tăng dần rồi ghi lại vị trí
bắt đầu của chúng.

\paragraph{Mảng hạng.}
\(\Rank[i]\) là hạng của \(\Suffix(i)\) trong thứ tự tăng dần của mọi hậu tố.
Một cách nhớ ngắn gọn là: mảng hậu tố trả lời ``ở hạng này là ai?'', còn mảng
hạng trả lời ``hậu tố này đứng thứ mấy?''. Hai mảng là nghịch đảo của nhau.
Nếu đã có một trong hai mảng, ta tính mảng còn lại trong \(\Oh(n)\).

\luosafigure{01}{0.88}{Quan hệ nghịch đảo giữa mảng hậu tố \(\SA\) và mảng hạng \(\Rank\) trên xâu \texttt{aabaaaab}.}

Khi cài đặt, để so sánh thuận tiện, có thể thêm vào cuối xâu một ký tự đặc
biệt chưa xuất hiện trước đó và nhỏ hơn mọi ký tự khác. Sau khi có mảng hạng,
ta so sánh hai hậu tố bất kỳ trong \(\Oh(1)\) bằng cách so sánh hạng. Nếu so
sánh trực tiếp từng ký tự thì trong trường hợp xấu nhất cần tới \(n\) phép so
sánh, nhưng chắc chắn không cần hơn.

\subsection{Thuật toán nhân đôi}

Ý tưởng chính của thuật toán nhân đôi là sắp xếp các đoạn con độ dài \(2^k\)
bắt đầu ở mọi vị trí và tính hạng của chúng. Ban đầu \(k=0\), tức mỗi đoạn
chỉ là một ký tự. Mỗi lần tăng \(k\) lên \(1\). Khi \(2^k>n\), mỗi đoạn độ
dài \(2^k\) tương đương với một hậu tố hoàn chỉnh; nếu mọi hạng đã khác nhau
thì ta đã có kết quả cuối cùng.

Trong mỗi vòng, hạng của đoạn độ dài \(2^k\) được biểu diễn bằng hai khóa:
hạng của nửa đầu độ dài \(2^{k-1}\) và hạng của nửa sau độ dài \(2^{k-1}\).
Sắp xếp hai khóa này bằng sắp xếp cơ số sẽ cho thứ tự mới. Với xâu
\texttt{aabaaaab}, các khóa \(x,y\) ở từng vòng lần lượt mô tả hai nửa của
đoạn đang xét; các hạng dần tách ra cho tới khi mọi hậu tố có hạng riêng.

\luosafigure{02}{0.9}{Quá trình nhân đôi trên xâu \texttt{aabaaaab}; hai khóa \(x,y\) biểu diễn hai nửa của đoạn đang sắp.}

\paragraph{Mã cài đặt.}
\begin{verbatim}
int wa[maxn],wb[maxn],wv[maxn],ws[maxn];
int cmp(int *r,int a,int b,int l)
{return r[a]==r[b]&&r[a+l]==r[b+l];}
void da(int *r,int *sa,int n,int m)
{
    int i,j,p,*x=wa,*y=wb,*t;
    for(i=0;i<m;i++) ws[i]=0;
    for(i=0;i<n;i++) ws[x[i]=r[i]]++;
    for(i=1;i<m;i++) ws[i]+=ws[i-1];
    for(i=n-1;i>=0;i--) sa[--ws[x[i]]]=i;
    for(j=1,p=1;p<n;j*=2,m=p)
    {
        for(p=0,i=n-j;i<n;i++) y[p++]=i;
        for(i=0;i<n;i++) if(sa[i]>=j) y[p++]=sa[i]-j;
        for(i=0;i<n;i++) wv[i]=x[y[i]];
        for(i=0;i<m;i++) ws[i]=0;
        for(i=0;i<n;i++) ws[wv[i]]++;
        for(i=1;i<m;i++) ws[i]+=ws[i-1];
        for(i=n-1;i>=0;i--) sa[--ws[wv[i]]]=y[i];
        for(t=x,x=y,y=t,p=1,x[sa[0]]=0,i=1;i<n;i++)
        x[sa[i]]=cmp(y,sa[i-1],sa[i],j)?p-1:p++;
    }
    return;
}
\end{verbatim}

Xâu cần sắp được đặt trong mảng \(r[0..n-1]\), mọi giá trị đều nhỏ hơn \(m\).
Để hàm thao tác thuận tiện, quy ước \(r[n-1]=0\), còn các \(r[i]\) khác đều
dương. Kết quả nằm trong \(sa[0..n-1]\).

Bước đầu tiên sắp xếp các xâu độ dài \(1\). Trong bài toán xâu thông thường,
giá trị ký tự không lớn, nên dùng sắp xếp cơ số:
\begin{verbatim}
for(i=0;i<m;i++) ws[i]=0;
for(i=0;i<n;i++) ws[x[i]=r[i]]++;
for(i=1;i<m;i++) ws[i]+=ws[i-1];
for(i=n-1;i>=0;i--) sa[--ws[x[i]]]=i;
\end{verbatim}
Ở đây \(x\) đóng vai trò mảng hạng tạm thời. Ta chưa cần nén thành hạng thật
sự, vì những bước sau chỉ dùng \(x\) để so sánh độ lớn của ký tự.

Trong mỗi vòng nhân đôi, sắp xếp cơ số trên hai khóa thường gồm hai lượt:
sắp theo khóa thứ hai rồi theo khóa thứ nhất. Kết quả sắp theo khóa thứ hai
có thể suy ra trực tiếp từ \(\SA\) của vòng trước, không cần sắp lại:
\begin{verbatim}
for(p=0,i=n-j;i<n;i++) y[p++]=i;
for(i=0;i<n;i++) if(sa[i]>=j) y[p++]=sa[i]-j;
\end{verbatim}
Biến \(j\) là độ dài nửa đoạn hiện tại, còn \(y\) lưu thứ tự đã sắp theo khóa
thứ hai. Sau đó chỉ cần sắp theo khóa thứ nhất:
\begin{verbatim}
for(i=0;i<n;i++) wv[i]=x[y[i]];
for(i=0;i<m;i++) ws[i]=0;
for(i=0;i<n;i++) ws[wv[i]]++;
for(i=1;i<m;i++) ws[i]+=ws[i-1];
for(i=n-1;i>=0;i--) sa[--ws[wv[i]]]=y[i];
\end{verbatim}

Khi đã có \(\SA\) mới, ta tính hạng mới. Có thể có nhiều đoạn nhận cùng hạng,
nên cần kiểm tra hai đoạn có giống hệt nhau hay không. Để tiết kiệm bộ nhớ,
mảng \(y\) được dùng lại làm mảng hạng cũ; hai con trỏ \(x,y\) được hoán đổi
thay vì sao chép cả mảng:
\begin{verbatim}
for(t=x,x=y,y=t,p=1,x[sa[0]]=0,i=1;i<n;i++)
x[sa[i]]=cmp(y,sa[i-1],sa[i],j)?p-1:p++;
\end{verbatim}
Hàm so sánh là
\begin{verbatim}
int cmp(int *r,int a,int b,int l)
{return r[a]==r[b]&&r[a+l]==r[b+l];}
\end{verbatim}
Quy ước \(r[n-1]=0\) giúp tránh kiểm tra biên: nếu \(r[a]=r[b]\) thì đoạn độ
dài \(l\) bắt đầu ở \(a\) và \(b\) chưa chứa ký tự kết thúc, nên truy cập
\(r[a+l]\) và \(r[b+l]\) vẫn hợp lệ. Sau vòng này, \(p\) là số hạng khác
nhau. Nếu \(p=n\), mọi đoạn đã phân biệt, các vòng sau không làm thay đổi thứ
tự nữa. Vì vậy điều kiện vòng lặp được viết:
\begin{verbatim}
for(j=1,p=1;p<n;j*=2,m=p) { ... }
\end{verbatim}

Mỗi vòng dùng sắp xếp cơ số trong \(\Oh(n)\); số vòng nhiều nhất là
\(\Oh(\log n)\). Độ phức tạp thời gian của thuật toán nhân đôi là
\(\Oh(n\log n)\), bộ nhớ là tuyến tính.

\subsection{Thuật toán DC3}

DC3 là thuật toán xây dựng mảng hậu tố tuyến tính theo tư tưởng chia để trị.
Thuật toán gồm ba bước.

\paragraph{Bước 1: sắp xếp các hậu tố có vị trí không chia hết cho 3.}
Chia hậu tố thành hai phần. Phần thứ nhất gồm các hậu tố bắt đầu ở vị trí
\(k\) với \(k\bmod 3\ne 0\), tức
\(\Suffix(1),\Suffix(2),\Suffix(4),\Suffix(5),\Suffix(7),\ldots\). Ta ghép
hai dãy vị trí \(1\bmod 3\) và \(2\bmod 3\). Nếu độ dài không phải bội của
3 thì thêm các số \(0\) ở cuối để đủ nhóm. Sau đó cứ ba ký tự được coi là một
bộ ba, sắp xếp cơ số các bộ ba này và nén mỗi bộ ba thành một ký tự mới. Mảng
hậu tố của xâu mới được tính đệ quy. Từ đó suy ra thứ tự của mọi hậu tố có
vị trí bắt đầu không chia hết cho \(3\). Xâu gốc phải kết thúc bằng một ký tự
nhỏ nhất và chưa từng xuất hiện để bảo đảm thứ tự suy ra là đúng.

\luosafigure{03}{0.82}{Bước đầu của DC3: lấy các hậu tố ở vị trí không chia hết cho \(3\), gom thành bộ ba, rồi tạo xâu mới để xử lý đệ quy.}

\paragraph{Bước 2: sắp xếp các hậu tố có vị trí chia hết cho 3.}
Mỗi hậu tố còn lại có dạng một ký tự đi kèm một hậu tố đã có hạng ở bước 1.
Do đó chỉ cần một lượt sắp xếp cơ số để có thứ tự của nhóm này.

\paragraph{Bước 3: trộn hai kết quả.}
Thao tác trộn giống bước trộn trong mergesort. Cần so sánh nhanh một hậu tố
thuộc nhóm \(0\bmod 3\) với một hậu tố thuộc nhóm còn lại. Có hai trường hợp:
\[
\begin{aligned}
  \Suffix(3i) &= r[3i] + \Suffix(3i+1),\\
  \Suffix(3j+1) &= r[3j+1] + \Suffix(3j+2),
\end{aligned}
\]
trong đó phần sau đã có hạng từ bước 1; và
\[
\begin{aligned}
  \Suffix(3i) &= r[3i] + r[3i+1] + \Suffix(3i+2),\\
  \Suffix(3j+2) &= r[3j+2] + r[3j+3] + \Suffix(3(j+1)+1).
\end{aligned}
\]
Lý do phải gộp theo bộ ba thay vì bộ đôi nằm ở đây: sau khi nhìn một hoặc hai
ký tự đầu, phần còn lại rơi vào nhóm đã được xếp hạng ở bước 1.

\paragraph{Mã cài đặt.}
\begin{verbatim}
#define F(x) ((x)/3+((x)%3==1?0:tb))
#define G(x) ((x)<tb?(x)*3+1:((x)-tb)*3+2)
int wa[maxn],wb[maxn],wv[maxn],ws[maxn];
int c0(int *r,int a,int b)
{return r[a]==r[b]&&r[a+1]==r[b+1]&&r[a+2]==r[b+2];}
int c12(int k,int *r,int a,int b)
{if(k==2) return r[a]<r[b]||r[a]==r[b]&&c12(1,r,a+1,b+1);
else return r[a]<r[b]||r[a]==r[b]&&wv[a+1]<wv[b+1];}
void sort(int *r,int *a,int *b,int n,int m)
{
    int i;
    for(i=0;i<n;i++) wv[i]=r[a[i]];
    for(i=0;i<m;i++) ws[i]=0;
    for(i=0;i<n;i++) ws[wv[i]]++;
    for(i=1;i<m;i++) ws[i]+=ws[i-1];
    for(i=n-1;i>=0;i--) b[--ws[wv[i]]]=a[i];
    return;
}
void dc3(int *r,int *sa,int n,int m)
{
    int i,j,*rn=r+n,*san=sa+n,ta=0,tb=(n+1)/3,tbc=0,p;
    r[n]=r[n+1]=0;
    for(i=0;i<n;i++) if(i%3!=0) wa[tbc++]=i;
    sort(r+2,wa,wb,tbc,m);
    sort(r+1,wb,wa,tbc,m);
    sort(r,wa,wb,tbc,m);
    for(p=1,rn[F(wb[0])]=0,i=1;i<tbc;i++)
    rn[F(wb[i])]=c0(r,wb[i-1],wb[i])?p-1:p++;
    if(p<tbc) dc3(rn,san,tbc,p);
    else for(i=0;i<tbc;i++) san[rn[i]]=i;
    for(i=0;i<tbc;i++) if(san[i]<tb) wb[ta++]=san[i]*3;
    if(n%3==1) wb[ta++]=n-1;
    sort(r,wb,wa,ta,m);
    for(i=0;i<tbc;i++) wv[wb[i]=G(san[i])]=i;
    for(i=0,j=0,p=0;i<ta && j<tbc;p++)
    sa[p]=c12(wb[j]%3,r,wa[i],wb[j])?wa[i++]:wb[j++];
    for(;i<ta;p++) sa[p]=wa[i++];
    for(;j<tbc;p++) sa[p]=wb[j++];
    return;
}
\end{verbatim}

Các tham số có ý nghĩa giống hàm nhân đôi. Khác biệt là \(r\) và \(sa\) nên
có kích thước khoảng \(3n\), để vùng nhớ phía sau được dùng cho xâu con đệ quy
mà không phải cấp phát lại. Trong hàm, \(rn=r+n\) lưu xâu mới ở bước 1, còn
\(san=sa+n\) lưu mảng hậu tố của xâu mới. Biến \(ta\) là số hậu tố ở vị trí
\(0\bmod 3\), \(tb=(n+1)/3\) là số hậu tố ở vị trí \(1\bmod 3\), và \(tbc\)
là tổng số hậu tố thuộc hai lớp \(1\bmod 3\) và \(2\bmod 3\).

Để tạo xâu mới, trước hết đặt hai ký tự canh ở cuối bằng \(0\), liệt kê các
vị trí không chia hết cho \(3\), rồi sắp xếp cơ số ba lần theo ba ký tự:
\begin{verbatim}
r[n]=r[n+1]=0;
for(i=0;i<n;i++) if(i%3!=0) wa[tbc++]=i;
sort(r+2,wa,wb,tbc,m);
sort(r+1,wb,wa,tbc,m);
sort(r,wa,wb,tbc,m);
\end{verbatim}
Hàm \texttt{sort} chỉ là một lượt sắp xếp cơ số theo khóa \(r[a[i]]\).

Sau khi các bộ ba đã được sắp, ta nén chúng thành các hạng mới:
\begin{verbatim}
for(p=1,rn[F(wb[0])]=0,i=1;i<tbc;i++)
rn[F(wb[i])]=c0(r,wb[i-1],wb[i])?p-1:p++;
\end{verbatim}
Ở đây \(F(x)\) đổi vị trí của hậu tố trong xâu gốc sang vị trí tương ứng
trong xâu mới, còn \(c0\) kiểm tra hai bộ ba có giống hệt nhau hay không.
Nếu mọi bộ ba đều khác nhau, ta đảo trực tiếp mảng hạng thành \(san\); nếu
không, gọi đệ quy:
\begin{verbatim}
if(p<tbc) dc3(rn,san,tbc,p);
else for(i=0;i<tbc;i++) san[rn[i]]=i;
\end{verbatim}

Tiếp theo là sắp nhóm vị trí chia hết cho \(3\). Thứ tự theo khóa thứ hai có
thể suy ra từ \(san\), rồi chỉ cần sắp theo ký tự đầu. Trường hợp \(n\bmod
3=1\) cần thêm hậu tố cuối \(n-1\), vì không có hậu tố \(n\) trong \(san\):
\begin{verbatim}
for(i=0;i<tbc;i++) if(san[i]<tb) wb[ta++]=san[i]*3;
if(n%3==1) wb[ta++]=n-1;
sort(r,wb,wa,ta,m);
for(i=0;i<tbc;i++) wv[wb[i]=G(san[i])]=i;
\end{verbatim}
Hàm \(G\) là ánh xạ ngược của \(F\), đổi vị trí trong xâu mới về vị trí trong
xâu gốc.

Cuối cùng trộn hai danh sách đã sắp:
\begin{verbatim}
for(i=0,j=0,p=0;i<ta && j<tbc;p++)
sa[p]=c12(wb[j]%3,r,wa[i],wb[j])?wa[i++]:wb[j++];
for(;i<ta;p++) sa[p]=wa[i++];
for(;j<tbc;p++) sa[p]=wb[j++];
\end{verbatim}
Hàm \(c12\) chính là phép so sánh hai trường hợp đã mô tả ở bước 3.

Nếu gọi độ phức tạp là \(f(n)\), các bước sắp xếp và trộn đều tuyến tính, còn
lời gọi đệ quy xử lý xâu dài không quá \(2n/3\):
\[
  f(n) = \Oh(n) + f(2n/3).
\]
Suy ra
\[
  f(n) \le c n + c(2n/3) + c(4n/9) + \cdots \le 3cn,
\]
nên DC3 là thuật toán \(\Oh(n)\).

\subsection{So sánh thuật toán nhân đôi và DC3}

\paragraph{Độ phức tạp thời gian.}
Nhân đôi chạy trong \(\Oh(n\log n)\), còn DC3 chạy trong \(\Oh(n)\). Tuy
nhiên hằng số của DC3 lớn hơn.

\paragraph{Độ phức tạp bộ nhớ.}
Cả hai đều dùng \(\Oh(n)\) bộ nhớ. Với cách cài đặt ở trên, nhân đôi cần tổng
kích thước mảng khoảng \(6n\), còn DC3 khoảng \(10n\).

\paragraph{Độ khó cài đặt.}
Bản nhân đôi ngắn hơn và dễ kiểm soát hơn; bản DC3 dài hơn một chút và có
nhiều ánh xạ chỉ số hơn.

\paragraph{Hiệu quả thực tế.}
Bài gốc thử nghiệm trên NOI-linux, CPU Pentium 4 2.80GHz, không tính thời
gian nhập xuất:
\[
\begin{array}{r|r|r}
N & \text{nhan doi (ms)} & \text{DC3 (ms)}\\
\hline
200000 & 192 & 140\\
300000 & 367 & 244\\
500000 & 750 & 499\\
1000000 & 1693 & 1248
\end{array}
\]
DC3 có ưu thế thực tế nhất định, còn nhân đôi dễ viết hơn. Khi làm bài, nên
chọn theo giới hạn dữ liệu và mức độ rủi ro cài đặt.

\section{Ứng dụng của mảng hậu tố}

Phần này trình bày các ứng dụng thường gặp. Mẫu chung là xây dựng mảng hậu tố
và mảng \(\Height\), rồi biến điều kiện về xâu con thành điều kiện trên các
hậu tố kề nhau hoặc trên các đoạn của mảng \(\Height\).

\subsection{Tiền tố chung dài nhất}

Định nghĩa
\[
  \Height[i] = \LCP(\Suffix(\SA[i-1]),\Suffix(\SA[i])).
\]
Với hai vị trí \(j,k\), giả sử \(\Rank[j]<\Rank[k]\), ta có tính chất quan
trọng:
\[
  \LCP(\Suffix(j),\Suffix(k))
  = \min_{\Rank[j] < t \le \Rank[k]} \Height[t].
\]
Trực giác là khi đi từ hậu tố này tới hậu tố kia trong thứ tự từ điển, mọi
cặp kề nhau trên đoạn phải cùng chia sẻ phần tiền tố chung; điểm yếu nhất
trên đoạn quyết định độ dài tiền tố chung của hai đầu mút. Chẳng hạn với xâu
\texttt{aabaaaab}, khi hỏi LCP của hai hậu tố \texttt{abaaaab} và
\texttt{aaab}, ta lấy giá trị nhỏ nhất của \(\Height\) trên đoạn giữa hai
hạng tương ứng.

\luosafigure{04}{0.86}{LCP của hai hậu tố bất kỳ bằng giá trị nhỏ nhất trên đoạn \(\Height\) giữa hai hạng của chúng.}

\paragraph{Tính \(\Height\) trong thời gian tuyến tính.}
Nếu tính \(\Height[2],\Height[3],\ldots,\Height[n]\) độc lập bằng cách so ký
tự, trường hợp xấu nhất là \(\Oh(n^2)\). Đặt
\[
  h[i]=\Height[\Rank[i]],
\]
tức \(h[i]\) là LCP của \(\Suffix(i)\) và hậu tố đứng ngay trước nó trong thứ
tự từ điển. Ta có
\[
  h[i]\ge h[i-1]-1.
\]
Thật vậy, giả sử \(\Suffix(k)\) đứng ngay trước \(\Suffix(i-1)\) và LCP của
chúng là \(h[i-1]\). Nếu \(h[i-1]>1\), sau khi bỏ ký tự đầu, \(\Suffix(k+1)\)
đứng trước \(\Suffix(i)\) và có LCP ít nhất \(h[i-1]-1\) với \(\Suffix(i)\).
Vì hậu tố đứng ngay trước \(\Suffix(i)\) là lựa chọn tốt nhất trong các hậu
tố đứng trước nó, \(h[i]\) không nhỏ hơn giá trị đó. Nếu \(h[i-1]\le 1\) thì
bất đẳng thức hiển nhiên.

Ta không cần lưu riêng mảng \(h\); chỉ cần duyệt \(i\) tăng dần và giữ biến
\(k\):
\begin{verbatim}
int rank[maxn],height[maxn];
void calheight(int *r,int *sa,int n)
{
    int i,j,k=0;
    for(i=1;i<=n;i++) rank[sa[i]]=i;
    for(i=0;i<n;height[rank[i++]]=k)
    for(k?k--:0,j=sa[rank[i]-1];r[i+k]==r[j+k];k++);
    return;
}
\end{verbatim}

\paragraph{Ví dụ 1: truy vấn tiền tố chung dài nhất.}
Cho một xâu và nhiều truy vấn, mỗi truy vấn hỏi LCP của hai hậu tố. Theo tính
chất trên, bài toán trở thành truy vấn giá trị nhỏ nhất trên một đoạn của
\(\Height\). Tiền xử lý \(\RMQ\) trong \(\Oh(n\log n)\) cho phép trả lời mỗi
truy vấn trong \(\Oh(1)\). Nếu dùng tiền xử lý \(\RMQ\) tuyến tính, toàn bộ
tiền xử lý cũng đạt \(\Oh(n)\).

\subsection{Các bài toán trên một xâu}

Với một xâu đơn, cách làm phổ biến là tính \(\SA\) và \(\Height\), sau đó
nhìn các xâu con như tiền tố của các hậu tố.

\subsubsection{Xâu con lặp}

\paragraph{Định nghĩa.}
Nếu xâu \(R\) xuất hiện ít nhất hai lần trong xâu \(L\), ta gọi \(R\) là một
xâu con lặp của \(L\).

\paragraph{Ví dụ 2: xâu con lặp dài nhất, được phép chồng nhau.}
Cho một xâu, tìm xâu con lặp dài nhất, hai lần xuất hiện có thể chồng nhau.
Độ dài cần tìm chính là giá trị lớn nhất của mảng \(\Height\). Lý do: LCP của
bất kỳ hai hậu tố đều là giá trị nhỏ nhất trên một đoạn của \(\Height\), nên
không thể lớn hơn phần tử lớn nhất của \(\Height\); ngược lại mỗi
\(\Height[i]\) là LCP của hai hậu tố kề nhau và tương ứng với một xâu lặp.
Thời gian sau khi có \(\Height\) là \(\Oh(n)\).

\paragraph{Ví dụ 3: xâu con lặp dài nhất không chồng nhau (PKU 1743).}
Nhị phân đáp án \(k\). Cần kiểm tra có tồn tại hai xâu con độ dài \(k\) bằng
nhau và không chồng nhau hay không. Dùng \(\Height\) để chia các hậu tố đã
sắp thành từng nhóm, sao cho các hậu tố trong cùng nhóm được nối bởi các giá
trị \(\Height\ge k\). Chỉ các hậu tố cùng nhóm mới có thể có LCP ít nhất
\(k\). Với mỗi nhóm, xét giá trị \(\SA\) nhỏ nhất và lớn nhất; nếu hiệu của
chúng ít nhất \(k\), hai lần xuất hiện không chồng nhau. Cách nhóm theo
\(\Height\) này rất thường gặp. Tổng thời gian \(\Oh(n\log n)\).

\luosafigure{05}{0.86}{Khi \(k=2\), xâu \texttt{aabaaaab} được chia thành các nhóm hậu tố nối bởi các giá trị \(\Height\ge k\).}

\paragraph{Ví dụ 4: xâu con lặp dài nhất xuất hiện ít nhất \(k\) lần (PKU
3261).}
Bài này tương tự ví dụ 3: nhị phân độ dài, chia nhóm bởi các vị trí có
\(\Height\ge k\), nhưng điều kiện kiểm tra là có nhóm chứa ít nhất \(k\) hậu
tố. Khi đó tồn tại một xâu con độ dài đang xét xuất hiện ít nhất \(k\) lần.
Tổng thời gian \(\Oh(n\log n)\).

\subsubsection{Số lượng xâu con}

\paragraph{Ví dụ 5: số xâu con khác nhau (SPOJ 694, SPOJ 705).}
Mọi xâu con đều là tiền tố của một hậu tố. Duyệt các hậu tố theo thứ tự
\[
  \Suffix(\SA[1]),\Suffix(\SA[2]),\ldots,\Suffix(\SA[n]).
\]
Khi thêm hậu tố \(\Suffix(\SA[k])\), nó có \(n-\SA[k]+1\) tiền tố. Trong số
đó, \(\Height[k]\) tiền tố đã xuất hiện trong các hậu tố trước, vì chúng là
tiền tố chung với hậu tố đứng trước trong thứ tự từ điển. Do đó đóng góp mới
là
\[
  n-\SA[k]+1-\Height[k].
\]
Cộng các giá trị này cho mọi \(k\) sẽ được số xâu con khác nhau trong
\(\Oh(n)\).

\subsubsection{Xâu con đối xứng}

\paragraph{Định nghĩa.}
Một xâu con là đối xứng nếu viết ngược nó vẫn thu được chính nó.

\paragraph{Ví dụ 6: xâu con đối xứng dài nhất (URAL 1297).}
Liệt kê tâm của xâu đối xứng. Cần xét riêng độ dài lẻ và độ dài chẵn. Cả hai
trường hợp đều quy về LCP giữa một hậu tố của xâu gốc và một hậu tố của xâu
đảo. Cách làm là ghép xâu gốc, một ký tự phân cách đặc biệt, rồi xâu đảo.
Sau đó xây dựng \(\SA\), \(\Height\), và cấu trúc \(\RMQ\) trên xâu ghép để
truy vấn LCP tương ứng với từng tâm. Khi phản chiếu hai nửa qua tâm, bài toán
trở thành hỏi LCP giữa một hậu tố trong xâu gốc và một hậu tố trong phần đảo.
Độ phức tạp là \(\Oh(n\log n)\), hoặc \(\Oh(n)\) nếu \(\RMQ\) được tiền xử lý
tuyến tính.

\luosafigure{06}{0.78}{Biến truy vấn xâu đối xứng quanh một tâm thành LCP trên xâu ghép từ xâu gốc và xâu đảo.}

\subsubsection{Xâu con lặp liên tiếp}

\paragraph{Định nghĩa.}
Nếu xâu \(L\) thu được bằng cách lặp một xâu \(S\) đúng \(R\) lần, ta gọi
\(L\) là xâu lặp liên tiếp, và \(R\) là số lần lặp.

\paragraph{Ví dụ 7: chu kỳ của cả xâu (PKU 2406).}
Cho xâu \(L\), biết nó là \(S\) lặp lại \(R\) lần; cần tìm \(R\) lớn nhất.
Liệt kê độ dài chu kỳ \(k\). Trước hết \(n\) phải chia hết cho \(k\). Sau đó
kiểm tra
\[
  \LCP(\Suffix(1),\Suffix(k+1))=n-k.
\]
Vì một đầu truy vấn luôn là \(\Suffix(1)\), không cần xây toàn bộ \(\RMQ\);
chỉ cần biết giá trị nhỏ nhất của \(\Height\) trên các đoạn nối tới
\(\Rank[1]\). Tổng thời gian \(\Oh(n)\).

\paragraph{Ví dụ 8: xâu con lặp liên tiếp nhiều lần nhất (SPOJ 687, PKU
3693).}
Liệt kê độ dài khối \(L\). Với mỗi \(L\), tìm số lần nhiều nhất mà một xâu
độ dài \(L\) có thể xuất hiện liên tiếp. Nếu có ít nhất hai khối liên tiếp,
đoạn đó phải phủ một cặp vị trí kề nhau trong dãy
\[
  0,L,2L,3L,\ldots.
\]
Với mỗi cặp \(L i\) và \(L(i+1)\), dùng LCP để mở rộng sang phải và dùng LCP
trên xâu đảo để mở rộng sang trái. Nếu tổng chiều dài khớp được là \(K\), số
lần lặp tại vùng đó là \(\lfloor K/L\rfloor+1\). Với mỗi \(L\), số cặp cần
xét là \(\Oh(n/L)\), nên tổng thời gian là
\[
  \Oh(n/1+n/2+\cdots+n/n)=\Oh(n\log n).
\]

\luosafigure{07}{0.7}{Với độ dài khối \(L\), xét hai mốc kề nhau \(Li\) và \(L(i+1)\), rồi mở rộng vùng khớp sang trái và sang phải.}

\subsection{Các bài toán trên hai xâu}

Với hai xâu, cách chuẩn là ghép chúng bằng một ký tự phân cách không xuất
hiện trong cả hai, xây dựng \(\SA\) và \(\Height\) trên xâu ghép, rồi chỉ
xét các hậu tố thuộc hai xâu khác nhau.

\subsubsection{Xâu con chung}

\paragraph{Định nghĩa.}
Nếu xâu \(L\) xuất hiện trong cả \(A\) và \(B\), ta gọi \(L\) là xâu con
chung của \(A\) và \(B\).

\paragraph{Ví dụ 9: xâu con chung dài nhất (PKU 2774, URAL 1517).}
Xâu con bất kỳ của một xâu là tiền tố của một hậu tố. Vì vậy xâu con chung
dài nhất của \(A\) và \(B\) là LCP lớn nhất giữa một hậu tố của \(A\) và một
hậu tố của \(B\). Sau khi ghép \(A\), ký tự phân cách, rồi \(B\), ta xét các
cặp hậu tố kề nhau trong \(\SA\). Với ví dụ \(A=\texttt{aaaba}\),
\(B=\texttt{abaa}\), hình dưới cho thấy chỉ các cặp kề nhau đến từ hai xâu
khác nhau mới được xét. Chỉ khi \(\SA[i-1]\) và \(\SA[i]\) thuộc
hai xâu gốc khác nhau thì \(\Height[i]\) mới là ứng viên hợp lệ. Giá trị lớn
nhất trong các ứng viên này là đáp án. Với \(|A|\) và \(|B|\) là độ dài hai
xâu, toàn bộ thuật toán chạy trong \(\Oh(|A|+|B|)\).

\luosafigure{08}{0.86}{Ghép \(A=\texttt{aaaba}\) và \(B=\texttt{abaa}\); chỉ các hậu tố kề nhau thuộc hai xâu khác nhau mới đóng góp cho xâu con chung dài nhất.}

\subsubsection{Số lượng xâu con chung}

\paragraph{Ví dụ 10: số xâu con chung độ dài ít nhất \(k\) (PKU 3415).}
Cho hai xâu \(A,B\), cần đếm số cặp xâu con chung có độ dài không nhỏ hơn
\(k\), cho phép các xâu con bằng nhau được tính nhiều lần. Ví dụ, với
\(A=\texttt{xx}\), \(B=\texttt{xx}\), \(k=1\), đáp án là \(5\). Với
\(A=\texttt{aababaa}\), \(B=\texttt{abaabaa}\), \(k=2\), đáp án là \(22\).

Ý tưởng là cộng, trên mọi cặp hậu tố một từ \(A\) và một từ \(B\), phần vượt
qua ngưỡng \(k\) của LCP. Sau khi ghép hai xâu và chia nhóm theo
\(\Height\), cần thống kê nhanh tổng LCP trong mỗi nhóm. Khi quét từ trái
sang phải, mỗi lần gặp một hậu tố của \(B\), ta tính nó tạo được bao nhiêu
xâu con chung độ dài ít nhất \(k\) với các hậu tố \(A\) đứng trước. Các hậu
tố \(A\) này được duy trì bằng một ngăn xếp đơn điệu theo giá trị
\(\Height\). Làm thêm một lượt đối xứng cho hậu tố \(A\) so với các hậu tố
\(B\) đứng trước. Chi tiết cài đặt là một bài tập về quét mảng \(\Height\) và
ngăn xếp đơn điệu.

\subsection{Các bài toán trên nhiều xâu}

Với nhiều xâu, ta ghép tất cả bằng các ký tự phân cách đôi một khác nhau và
không xuất hiện trong dữ liệu. Sau đó xây dựng \(\SA\), \(\Height\), rồi dùng
\(\Height\) để chia nhóm. Nhiều bài cần nhị phân đáp án.

\paragraph{Ví dụ 11: xâu dài nhất xuất hiện trong ít nhất \(k\) xâu (PKU
3294).}
Ghép \(n\) xâu bằng các ký tự phân cách riêng, xây dựng mảng hậu tố. Nhị phân
độ dài \(L\); với mỗi \(L\), chia nhóm bởi các vị trí có \(\Height\ge L\) và
kiểm tra có nhóm nào chứa hậu tố đến từ ít nhất \(k\) xâu gốc hay không. Độ
phức tạp \(\Oh(n\log n)\), với \(n\) là tổng độ dài.

\paragraph{Ví dụ 12: xâu dài nhất xuất hiện ít nhất hai lần không chồng nhau
trong mỗi xâu (SPOJ 220).}
Cách làm gần giống ví dụ 11. Với một độ dài đang xét, chia nhóm theo
\(\Height\). Một nhóm hợp lệ nếu trong mỗi xâu gốc có ít nhất hai hậu tố thuộc
nhóm, và trong từng xâu đó, hiệu giữa vị trí bắt đầu lớn nhất và nhỏ nhất
không nhỏ hơn độ dài đang xét. Điều kiện hiệu vị trí bảo đảm hai lần xuất
hiện không chồng nhau; nếu đề không yêu cầu không chồng nhau thì bỏ kiểm tra
này. Độ phức tạp \(\Oh(n\log n)\).

\paragraph{Ví dụ 13: xâu dài nhất xuất hiện trực tiếp hoặc sau khi đảo trong
mỗi xâu (PKU 3294).}
Điểm khác biệt là mỗi xâu có thể xuất hiện ở dạng đảo. Ta chỉ cần thêm bản
đảo của từng xâu vào chuỗi ghép, vẫn dùng các ký tự phân cách đôi một khác
nhau. Sau đó nhị phân đáp án và chia nhóm theo \(\Height\) như trước. Khi
kiểm tra một nhóm, với mỗi xâu gốc chỉ cần có hậu tố thuộc bản thuận hoặc bản
đảo. Độ phức tạp vẫn là \(\Oh(n\log n)\).

\section{Kết luận}

Mảng hậu tố là một cấu trúc dữ liệu rất tốt cho xử lý xâu và là một công cụ
mạnh trong nhiều dạng bài. Người học nên nắm vững cách xây dựng mảng hậu tố,
mảng hạng và mảng \(\Height\), đồng thời luyện cách chuyển điều kiện về xâu
con thành điều kiện trên các hậu tố liên tiếp hoặc trên các đoạn của
\(\Height\). Các kỹ thuật như chia nhóm theo ngưỡng \(\Height\), nhị phân đáp
án, ghép nhiều xâu bằng ký tự phân cách và dùng \(\RMQ\) cho LCP xuất hiện
lặp lại trong rất nhiều bài toán.

\section*{Tài liệu tham khảo}

\begin{enumerate}
\item Liu Rujia, \textit{Nghệ thuật thuật toán và thi Tin học}, Nhà xuất bản
Đại học Thanh Hoa, Bắc Kinh, 2004.
\item Xu Zhilei, luận văn đội tuyển quốc gia IOI 2004, \textit{Mảng hậu tố}.
\item Zhou Yuan, bài tập đội tuyển quốc gia Tin học 2005, \textit{Thuật toán
sắp xếp hậu tố tuyến tính}.
\end{enumerate}

\section*{Lời cảm ơn}

Tác giả cảm ơn CCF đã cung cấp một diễn đàn để trao đổi với mọi người; cảm
ơn huấn luyện viên Tang Wenbin của Đại học Thanh Hoa; cảm ơn thầy Zhang
Xuedong đã giúp đỡ trong quá trình viết bài; cảm ơn thầy Lin Shenghua của
Trường Trung học số 2 Quảng Châu đã hướng dẫn và gợi mở; cảm ơn Fang Zhanpeng
của Trường Trung học số 1 Trung Sơn và Jiang Biye của Trường Trung học Kỷ
niệm Tôn Trung Sơn đã cung cấp tư liệu về mảng hậu tố.

\end{document}
